\documentclass[border=4pt]{standalone}
\usepackage{amsmath,amssymb,mathtools,xcolor,varwidth}
\begin{document}
\begin{varwidth}{28em}
\large\bfseries
Kepler had noticed a haunting regularity in the heavens: the farther a planet      orbits from the Sun, the longer its year — and the relationship is exact. Squaring      the periodic time $T$ always tracks the cube of the greater axis $a$,      so that $T^2$ $\propto$ $a^3$. Equivalently, the period      grows as the $3/2$ power — what Newton calls the sesquiplicate ratio —      of the axis. What Kepler found only by patient observation, Newton now proves from      the inverse-square law: "the periodic times in ellipses are in the sesquiplicate ratio of their greater axes."
\end{varwidth}
\end{document}
